\def\stat{stepchenkov}



\def\tit{ФОРМАЛИЗАЦИЯ СИНТЕЗА САМОСИНХРОННЫХ~СЧЕТЧИКОВ}



\def\titkol{Формализация синтеза самосинхронных счетчиков}



\def\aut{Ю.\,А.~Степченков$^1$, Ю.\,Г.~Дьяченко$^2$, Н.\,В.~Морозов$^3$, 

Д.\,Ю.~Степченков$^4$, Д.\,Ю.~Дьяченко$^5$}



\def\autkol{Ю.\,А.~Степченков, Ю.\,Г.~Дьяченко, Н.\,В.~Морозов и~др.} 

%Д.\,Ю.~Степченков$^4$, Д.\,Ю.~Дьяченко$^5$}



\titel{\tit}{\aut}{\autkol}{\titkol}



\index{Степченков Ю.\,А.}

\index{Дьяченко Ю.\,Г.}

\index{Морозов Н.\,В.}

\index{Степченков Д.\,Ю.}

\index{Дьяченко Д.\,Ю.}

\index{Stepchenkov Yu.\,A.}

\index{Diachenko Yu.\,G.}

\index{Morozov N.\,V.}

\index{Stepchenkov D.\,Yu.} 

\index{Diachenko D.\,Yu.}



%{\renewcommand{\thefootnote}{\fnsymbol{footnote}}

%\footnotetext[1]

%{Исследование выполнено за счет гранта Российского научного фонда №\,22-79-10053 ({\sf https://rscf.ru/project/22-79-10053}).}} 



\renewcommand{\thefootnote}{\arabic{footnote}}

\footnotetext[1]{Федеральный исследовательский центр <<Информатика и~управ\-ле\-ние>> Российской академии наук, 

\mbox{YStepchenkov@ipiran.ru}}

\footnotetext[2]{Федеральный исследовательский центр <<Информатика и~управ\-ле\-ние>> Российской академии наук,  \mbox{diaura@mail.ru}}

\footnotetext[3]{Федеральный исследовательский центр <<Информатика и~управ\-ле\-ние>> Российской академии наук, \mbox{NMorozov@ipiran.ru}}

\footnotetext[4]{Федеральный исследовательский центр <<Информатика и~управление>> 

   Российской академии наук,  \mbox{stepchenkov@mail.ru}}

\footnotetext[5]{Федеральный исследовательский центр 

<<Информатика и~управление>> Российской академии наук, \mbox{diaden87@gmail.com}}





 \label{st\stat}



 \thispagestyle{headings}

 

\vspace*{-12pt}



  

  



\Abst{Самосинхронные (СС) схемы обладают высокой на\-деж\-ностью. 

Они гарантируют обнаружение и~локализацию любых константных 

неисправностей и~демонстрируют высокий уровень 

сбое\-устой\-чи\-вости. Однако проектирование СС-схем более трудоемко 

в~сравнении с~синхронными аналогами из-за не\-об\-хо\-ди\-мости по\-стро\-ения 

дополнительной индикаторной под\-схе\-мы и~соблюдения принципов 

функционирования схем, относящихся к~классу СС-схем. Автоматическое 

преобразование исходной схемы, описанной аппаратно как синхронная схема, 

в~СС-ре\-а\-ли\-за\-цию обеспечивается с~помощью процедуры 

формализованной десинхронизации. Но при синтезе последовательностных 

СС-устройств, в~том чис\-ле СС-счет\-чи\-ков, формальная десинхронизация 

приводит к~чрезмерной аппаратной из\-бы\-точ\-ности и,~как следствие, к~их 

низкому быст\-ро\-дей\-ст\-вию. Статья обосновывает подход к~синтезу  

СС-счет\-чи\-ков, ба\-зи\-ру\-ющий\-ся на формализации эвристических методов их 

по\-стро\-ения и~обес\-пе\-чи\-ва\-ющий гарантированное получение действительно 

СС-ре\-а\-ли\-за\-ции, функ\-ци\-о\-ни\-ру\-ющей в~пол\-ном соответствии с~исходным 

описанием и~об\-ла\-да\-ющей близ\-ки\-ми к~оптимальным по\-тре\-би\-тель\-ски\-ми 

характеристиками.}



\KW{автоматизированный синтез; самосинхронная схема; счетчик; 

десинхронизация; предуста\-нов\-ка; индикация}



\DOI{10.14357\!/\!08696527240205}{KDIEOJ}  



\vskip 10pt plus 9pt minus 6pt



 \thispagestyle{headings}

 

 \vspace*{-22pt}



\section{Введение}



\vspace*{-4pt}

  

  Успешное развитие современной экономики невозможно без широкого 

внед\-ре\-ния циф\-ро\-вой микроэлектроники во все области де\-я\-тель\-ности 

человека. Циф\-ро\-ви\-за\-ция жиз\-не\-де\-я\-тель\-ности человека по\-рож\-да\-ет сильную 

за\-ви\-си\-мость экономики и~общества от циф\-ро\-вой техники, от на\-деж\-ности, 

без\-от\-каз\-ности аппаратуры на интегральных мик\-ро\-схе\-мах. Современная 

мик\-ро\-элект\-ро\-ни\-ка реализуется в~основном на принципах синхронного 

управ\-ле\-ния процессами, ба\-зи\-ру\-юще\-го\-ся на глобальном <<дереве>> 

синхронизации~[1]. С~увеличением слож\-ности мик\-ро\-схем и~снижением 

норм топологического проектирования рас\-тет площадь крис\-тал\-лов 

синхронных интегральных мик\-ро\-схем и,~как следствие, объем аппаратных 

за\-трат и~энер\-го\-по\-треб\-ле\-ния, приходящихся на так\-то\-вое <<дерево>>. 

Разработчики синхронных схем вы\-нуж\-де\-ны за\-кла\-ды\-вать в~про\-ек\-ти\-ру\-емые 

синхронные схемы <<запас проч\-ности>> в~расчете на наихудший случай 

и~из\-бы\-точ\-ность аппаратуры, обес\-пе\-чи\-ва\-ющий ее ра\-бо\-то\-спо\-соб\-ность 

в~из\-ме\-ня\-ющих\-ся условиях эксплуатации и~в~процессе старения~[2].

  

  Альтернативой синхронным схемам служат асинхронные схемы, 

и~в~част\-ности их подкласс~--- СС-схе\-мы. В~них нет 

глобальной синхронизации. Они функционируют на основе  

за\-прос-от\-вет\-ной дис\-цип\-ли\-ны~[3--5]. Самосинхронные схе\-мы свободны от 

пе\-ре\-чис\-лен\-ных недостатков синхронных аналогов, гарантируют на\-де\-жную 

работу при любых предельных условиях эксплуатации и~характеризуются 

повышенной устой\-чи\-востью к~логическим сбоям~[6, 7]. Тем не менее 

по\-дав\-ля\-ющее большинство практических циф\-ро\-вых схем относятся 

к~синхронному классу. Развитию СС-схем препятствует большая слож\-ность 

и~не\-стан\-дарт\-ность процесса их руч\-ной раз\-ра\-бот\-ки и~отсутствие программных 

средств, обес\-пе\-чи\-ва\-ющих автоматизацию всего марш\-ру\-та проектирования 

\mbox{СС-ап}\-па\-ра\-ту\-ры и~при этом не тре\-бу\-ющих высокой квалификации от 

раз\-ра\-бот\-чика.

  

  Современные системы автоматизированного проектирования (САПР) 

циф\-ро\-вых устройств, ис\-поль\-зу\-емые в~промышленных мас\-шта\-бах, 

обеспечивают эффективное проектирование схем синхронного типа. Они 

поддерживаются и~развиваются за\-ру\-беж\-ны\-ми фирмами, например Cadence, 

Synopsys, Mentor Graphics и~др. Использование синхронных САПР для 

разработки СС-схем приводит к~неутешительным результатам: 

характеристики полученных СС-схем многократно уступают синхронным 

аналогам~[8, 9]. Промышленных САПР СС-схем пока нет. Сейчас за 

рубежом есть несколько программных средств разработки асинхронных 

схем, например Balsa~[10]. Они служат инструментом синтеза асинхронной 

схемы, используя в~качестве исходного задания ее описание на 

поведенческом или алгоритмическом уровне. При этом они используют 

форму исходного описания син\-те\-зи\-ру\-емой схемы, специфичную именно для 

СС-схем: специальный язык, сигнальные графы и~т.\,д. Такой подход 

за\-труд\-ня\-ет работу проектировщиков, но не гарантирует получение 

корректного с~точ\-ки зрения СС-дис\-цип\-ли\-ны результата~[11].

  

  Более простым и~интуитивно понятным представляется подход к~синтезу 

\mbox{СС-схем}, преобразующий исходное описание схемы, заданное в~синхронном 

виде, в~асинхронное с~по\-мощью метода десинхронизации~[12--14].  

С~по\-мощью анализа связ\-ности устройств в~со\-ста\-ве схемы, формальных 

правил замены синхронных элементов их СС-ана\-ло\-га\-ми и~замены 

глобальной синхронизации со\-во\-куп\-ностью локальных сигналов  

управ\-ле\-ния за\-прос-от\-вет\-ным взаимодействием этот метод позволяет 

достаточно прос\-то получить асинхронную схему. Однако синтезированные 

таким путем схемы используют для формирования сигналов управ\-ле\-ния 

элементы за\-держ\-ки вмес\-то проверки и~под\-тверж\-де\-ния завершения 

переключения схемы в~текущую фазу работы (индицирования~[15]) и~из-за 

этого теряют ряд преимуществ полноценных СС-схем, в~част\-ности 

не\-за\-ви\-си\-мость работы от задержек элементов. 

  

  Отечественные программные средства проектирования СС-схем касаются 

лишь вопросов диагностирования~[16] и~анализа на са\-мо\-синхрон\-ность~[17]. 

Они обеспечивают проверку кор\-рект\-ности СС-ре\-а\-ли\-за\-ции синтезированной 

схемы, но не способны об\-лег\-чить и~ускорить ее разработку.

  

  Кроме того, прямое формальное преобразование синхронного описания 

  в~СС-схе\-му приводит к~большой из\-бы\-точ\-ности последовательностных  

СС-устройств, например двоичных счетчиков. Цель \mbox{статьи}~--- оценка 

воз\-мож\-ности применения име\-ющих\-ся программных средств синтеза 

циф\-ро\-вых устройств для синтеза двоичных \mbox{СС-счет}\-чи\-ков и~разработка 

формализованного метода синтеза \mbox{СС-счет}\-чи\-ков, обеспечивающего 

получение корректных схем с~приемлемыми характеристиками. Научная 

новизна опи\-сы\-ва\-емо\-го подхода заключается в~разработке алгоритма 

формализованного преобразования исходного Verilog-опи\-са\-ния двоичного 

синхронного счетчика в~СС-ре\-а\-ли\-за\-цию, функ\-ци\-о\-ни\-ру\-ющую в~полном 

соответствии с~исходным описанием с~учетом специфики работы СС-схем.



\begin{figure}[b] %fig1

\vspace*{-6pt}

  \begin{center}

  {\small \begin{tabular}{|ll|}

  \hline

  always@(posedge T or posedge res) begin &\\

  \hspace*{8mm}if(res) &/\!\!/  {\footnotesize Начальный сброс по условию} res\;=\;1\\

  \hspace*{11mm}State\;$<=$\;4$^\prime$h0;&\\

  \hspace*{8mm}else &/\!\!/ {\footnotesize Инкрементация счетчика}\\

  \hspace*{11mm}State\;$<=$\;State\;+\;1; &\\

  end&\\

  \hline

  \end{tabular}

  }

  \end{center}

  

  \vspace*{-6pt}

  

    \Caption{Verilog-описание двоичного счетчика}

  \end{figure}

  



\section{Автоматизация логического синтеза самосинхронных счетчиков}

  

  Формализованный синтез СС-счет\-чи\-ков лежит в~рус\-ле об\-щей 

методологии автоматизации проектирования СС-схем~[18]. Пусть исходное 

описание син\-те\-зи\-ру\-емой схемы пред\-став\-ле\-но на языке Verilog на 

поведенческом уровне с~использованием блоков always. Распознать 

синхронный счетчик в~исходном Verilog-опи\-са\-нии достаточно просто: 

в~блоке always, опи\-сы\-ва\-ющем счетчик, присутствует оператор инкремента 

или декремента векторной переменной по фронту (переднему или зад\-не\-му) 

некоторого сигнала, чаще всего глобального тактового сигнала. 

Дополнительные операторы в~теле блока always, описывающего счет\-чик, 

специфицируют его предуста\-нов\-ку по внеш\-не\-му сигналу разрешения или по 

заданному со\-сто\-янию самого счет\-чика. 

  

  Рисунок~1 демонстрирует пример описания двоичного счетчика 

с~асинхронным сбросом на языке Verilog. Состояние State инкрементируется 

при по\-ступ\-ле\-нии положительного фронта сигнала~T и~пассивном уровне 

входа сброса res, что соответствует сум\-ми\-ру\-юще\-му счетчику. Поскольку 

в~теле блока always нет дополнительных условий форматирования со\-сто\-яния 

State (восьмеричного, десятичного кодирования), это двоичный счетчик.

  

  

  Результат синтеза двоичного счетчика по его исходному поведенческому 

описанию зависит от уровня формализации алгоритма преобразования блока 

always в~схему на элементах биб\-лио\-те\-ки стандартных ячеек. Например, 

под\-сис\-те\-ма Yosys~[19] реализует любой двоичный счетчик в~виде ре\-гист\-ра 

хранения на двухтактных триггерах, со\-сто\-яние которого изменяется 

в~соответствии с~заданным типом счета (сум\-ми\-ру\-ющим, вы\-чи\-та\-ющим, 

реверсивным) с~по\-мощью параллельной записи в~разряды ре\-гист\-ра 

информации, подготовленной комбинационной схемой. Рисунок~2 

ил\-люст\-ри\-ру\-ет результат синтеза под\-сис\-те\-мой Yosys синхронного 

счетчика, заданного Verilog-ко\-до\-м на рис.~1.

  

  Такой подход к~синтезу счетчиков обеспечивает по\-стро\-ение в~едином 

стиле всех разновидностей счетчиков, с~любым типом предуста\-нов\-ки 

и~на\-прав\-ле\-ния счета. Информационный вход $i$-го разряда регистра 

вы\-чис\-ля\-ет\-ся как функция

  \begin{equation}

  X_i= \left( \mathop{\wedge}\limits_{k=0}^{i-1} \mathrm{State}\,[k]\right) \oplus \mathrm{State}\,[i]\,,

  \label{e1-step}

  \end{equation}

где символ <<$\wedge$>> означает логическое произведение. При $i\hm = 0$ 

функция~(1) упрощается до вида: 

$$

X_0= \sim \mathrm{State}\,[0]\,. 

$$





\begin{figure} %fig2

      \vspace*{1pt}

  \begin{center}

 \mbox{%

 \epsfxsize=125.412mm 

\epsfbox{ste-2.eps}

 }

\end{center}

\vspace*{-9pt}

\Caption{Результат синтеза Yosys суммирующего счетчика с~асинхронным сбросом}

\end{figure}



  Схема на рис.~2 логически обоснована, так как обеспечивает 

максимальное быст\-ро\-дей\-ст\-вие в~синхронной схемотехнике за счет 

избыточности аппаратных за\-трат, которая, однако, не всегда оказывается 

оправ\-дан\-ной. Суммарная аппаратная слож\-ность схемы счетчика на рис.~2 

рав\-на~180~транзисторам при изготовлении в~технологии КМДП 

(комплементарный ме\-талл\,--\,ди\-элект\-рик\,--\,по\-лу\-про\-вод\-ник). 

Превратить такую схему в~СС-ана\-лог можно путем прямого преобразования 

с~по\-мощью сле\-ду\-ющей процедуры, ре\-а\-ли\-зу\-ющей методы, изложенные 

в~работах~[18, 20]:

  \begin{itemize}

\item дуализировать комбинационную часть, используя традиционное для 

комбинационных СС-схем парафазное со спейсером кодирование всех 

информационных сигналов;

\item подставить вместо синхронных двухтактных триггеров, обозначенных 

именем FDC на рис.~2, подходящий двухтактный СС-триг\-гер с~входом 

управ\-ле\-ния и~парафазным со спейсером информационным входом; 

\item ввести преобразователь бифазных выходов СС-триг\-ге\-ров 

в~парафазные со спейсером сигналы, поскольку комбинационная часть 

СС-счет\-чи\-ка работает именно с~такими сигналами;

\item добавить индикаторную подсхему, учи\-ты\-ва\-ющую не\-об\-хо\-ди\-мость 

дополнительной индикации входа управ\-ле\-ния либо в~каж\-дом триггере, 

либо групповым методом из-за парафазного со спейсером типа 

информационного входа разрядов СС-счет\-чика. 

\end{itemize}



  В результате получается СС-схе\-ма, показанная на рис.~3 в~базисе 

биб\-лио\-те\-ки стандартных ячеек для проектирования СС-схем~[21]. Схема 

использованного в~ней СС-триг\-ге\-ра R1RE24 пред\-став\-ле\-на на рис.~4. Здесь 

R, S~--- информационный парафазный вход с~нулевым спейсером; RT~--- 

вход асинхронного сброса; Е~--- вход управ\-ле\-ния, аналог тактового сигнала 

в~синхронных триггерах; Q, QB~--- бифазный информационный выход; I~--- 

индикаторный выход. Суммарная аппаратная слож\-ность схемы на рис.~3 

рав\-на 376~КМДП-тран\-зис\-то\-рам.



\begin{figure} %fig3

 \vspace*{1pt}

  \begin{center}

 \mbox{%

 \epsfxsize=120.45mm  %128mm

\epsfbox{ste-3.eps}

 }

\end{center}

%\vspace*{-12pt}

%\Caption{Результат преобразования параллельного счетчика в~СС-реализацию}

\vspace*{-7.6pt}

\end{figure}



  В то же время максимальное быстродействие за счет из\-бы\-точ\-ности,  

во-пер\-вых, далеко не всегда востребовано, а~во-вто\-рых,  

в~СС-ва\-ри\-ан\-те оказывается проб\-ле\-ма\-тич\-ным вследствие чрезмерного 

услож\-не\-ния схемы. Для решения вопроса це\-ле\-со\-об\-раз\-ности 

использования параллельного СС-счет\-чи\-ка сравним его 

с~последовательным вариантом СС-счет\-чика.

  

\section{Сравнение вариантов самосинхронных счетчиков}

  

  Последовательный СС-счет\-чик также синтезируется на основе 

последовательного синхронного аналога. Рисунок~5 демонстрирует 

реализацию последовательного синхронного счетчика с~асинхронным 

сбро\-сом. Его аппаратная слож\-ность~--- 116~КМДП-тран\-зис\-то\-ров.



\setcounter{figure}{3}

\begin{figure} %fig4

\vspace*{1pt}

  \begin{center}

 \mbox{%

 \epsfxsize=81.597mm  

\epsfbox{ste-4.eps}

 }

\end{center}

\vspace*{-9pt}

\Caption{Схема триггера R1RE24}

\end{figure}



\begin{figure} %fig5

\vspace*{1pt}

  \begin{center}

 \mbox{%

 \epsfxsize=127.853mm  %128mm

\epsfbox{ste-5.eps}

 }

\end{center}

\vspace*{-9pt}

\Caption{Последовательный синхронный двоичный счетчик}

\end{figure}





\pagebreak



\



\vspace*{-12pt}





\begin{floatingfigure}{60mm} %fig6

\vspace*{-6pt}

  \begin{center}

 \mbox{%

 \epsfxsize=52.975mm  %128mm

\epsfbox{ste-6.eps}

 }

\end{center}

\vspace*{-9pt}

\Caption{Схема триггера С1С}

\end{floatingfigure}



  Она превращается в~СС-реализацию с~асинхронным сбросом путем 

замены триггеров, обозначенных именем FDC на рис.~5, триггером~С1С из  

биб\-лио\-те\-ки~[21] и~добавления индикаторной подсхемы. Схема 

триг\-ге\-ра~С1С изоб\-ра\-же\-на на рис.~6. Здесь T~--- счетный вход с~нулевым 

спейсером; С~--- вход СС сброса; Q, QB~--- бифазный 

информационный выход; I~--- индикаторный выход; OT~--- счетный выход; 

U, UB~--- внут\-рен\-ний бифазный сигнал, выход би\-ста\-биль\-ной ячейки. 

  





  Результирующая схема четырехразрядного последовательного  

СС-счет\-чи\-ка показана на рис.~7. Ее аппаратная слож\-ность~---  

154~КМДП-тран\-зис\-то\-ра~--- почти в~2,5~раза ниже, чем в~параллельном 

СС-счетчике на рис.~3.

  





  Сравнение вариантов четырехразрядного СС-счет\-чи\-ка, пред\-став\-лен\-ных 

на рис.~3 и~7, показывает, что в~СС-ре\-а\-ли\-за\-ции, в~отличие от 

синхронной схе\-мо\-тех\-ни\-ки, параллельный счетчик не самый быст\-рый. 

Моделирование на транзисторном уров\-не (программа Spectre, Cadence)  

в~65-на\-но\-мет\-ро\-вой КМДП-тех\-но\-ло\-гии, результаты которого иллюстрирует рис.~8, 

выявило сле\-ду\-ющие различия в~поведении вариантов СС-счет\-чика:\\[-13pt]



\begin{figure} %fig7

\vspace*{1pt}

  \begin{center}

 \mbox{%

 \epsfxsize=127.106mm  

\epsfbox{ste-7.eps}

 }

\end{center}

\vspace*{-13pt}

\Caption{Самосинхронная реализация последовательного счетчика с~асинхронным \mbox{сбросом} }

%\end{figure}

%\begin{figure}[b] %fig8

\vspace*{5pt}

  \begin{center}

 \mbox{%

 \epsfxsize=114.56mm  

\epsfbox{ste-8.eps}

 }

\end{center}

\vspace*{-13pt}

\Caption{Моделирование вариантов последовательного (выходы~Z$^*$ и~IndZ) 

и~параллельного (выходы Y$^*$ и~IndY) четырехразрядного СС-счет\-чи\-ка с~асинхронным 

сбросом}

\vspace*{-6pt}

\end{figure}



\noindent

\begin{enumerate}[(1)]

\item  в~параллельном варианте (см.\ рис.~3) времена пребывания СС-счет\-чи\-ка 

в~рабочей и~спейсерной фазах примерно рав\-ны и~не зависят от текущего 

со\-сто\-яния счетчика, в~то время как в~последовательном счетчике (см.\ рис.~7) 

дли\-тель\-ности фаз варь\-и\-ру\-ют\-ся в~за\-ви\-си\-мости от чис\-ла разрядов, 

пе\-ре\-клю\-ча\-ющих\-ся в~процессе перехода из одного со\-сто\-яния счетчика 

в~другое;\\[-13pt]

\item  в~автономном режиме, когда выход общего индикатора счетчика 

замкнут на его счетный вход, среднее быст\-ро\-дей\-ст\-вие СС-счет\-чи\-ка 

(величина, обрат\-ная по отношению к~времени прохода счетчика по всем 

со\-сто\-яни\-ям) в~по\-сле\-до\-ва\-тель\-ном варианте оказывается в~1,7~раза выше, 

чем в~параллельном счетчике.

\end{enumerate}







  Моделирование двух аналогичных вариантов реализации 

восьмиразрядного СС-счет\-чи\-ка демонстрирует, что с~увеличением 

раз\-ряд\-ности преимущество последовательного счетчика над параллельным 

растет. На рис.~9 показаны диаграммы переключения четырех старших 

разрядов восьмиразрядных последовательного и~параллельного  

СС-счет\-чи\-ков. Последовательный счетчик работает поч\-ти в~2~раза 

быст\-рее. 

  



\begin{figure} %fig9

\vspace*{1pt}

  \begin{center}

 \mbox{%

 \epsfxsize=114.24mm  

\epsfbox{ste-9.eps}

 }

\end{center}

\vspace*{-9pt}

\Caption{Моделирование вариантов восьмиразрядного СС-счет\-чи\-ка с~асинхронным 

сбросом}

\vspace*{-6pt}

\end{figure}



  Причина преимущества последовательного счетчика в~быст\-ро\-дей\-ст\-вии 

заключается в~более прос\-той схеме и~индикации в~сравнении с~параллельным 

СС-счет\-чи\-ком. В~параллельном варианте функция~(1) информационного 

входа $i$-го разряда счетчика требует декомпозиции час\-ти 

$\mathop{\wedge}_{k=0}^{i-1} \mathrm{State}\,[k]$ при $i \hm> 4$, поэтому 

с~увеличением раз\-ряд\-ности комбинационная часть параллельного счетчика 

становится все более декомпозированной, увеличиваются слож\-ность 

и~за\-держ\-ка индикаторной под\-схемы. 

  

  Дополнительное преимущество последовательного варианта  

СС-счет\-чи\-ка~--- его прос\-тая мас\-шта\-би\-ру\-емость, что облегчает его 

па\-ра\-мет\-ри\-за\-цию при автоматизированном синтезе. Поэтому для  

СС-ре\-а\-ли\-за\-ции целесообразно использовать последовательный  

счет\-чик.

  

\section{Формализация синтеза самосинхронных счетчиков}

  

  Самосинхронный счетчик целесообразно синтезировать на основе библиотеки 

параметризованных СС-счет\-чи\-ков. Пример параметризованного  

Verilog-опи\-са\-ния модуля СС-счет\-чи\-ка со сбросом представлен на 

рис.~10. 







  Параметр WIDTH задает разрядность счетчика и~настраивается под 

тре\-бу\-емое значение в~каж\-дом конкретном случае в~процессе синтеза, причем 

характер сброса~--- асинхронный или самосинхронный~--- не влияет на 

внут\-рен\-нюю реализацию счетчика. Соответствующая специфика с~успехом 

определяется окру\-же\-ни\-ем счетчика. Выходы модуля счетчика поз\-во\-ля\-ют 

организовать СС-сброс в~окру\-же\-нии счет\-чи\-ка. Раз\-ряд\-ность счетчика WIDTH 

извлекается из тела блока always или из объявления со\-от\-вет\-ст\-ву\-ющей 

переменной в~исходном Verilog-опи\-са\-нии. Индикаторная под\-схе\-ма 

выполняет функцию логического~<<И>> от со\-во\-куп\-ности выходов~I всех 

разрядов и~выхода OT по\-след\-не\-го (стар\-ше\-го) разряда:

  $$

  \mathrm{Ind}= \mathrm{OT}_{N} \mathop{\wedge}\left( \mathop{\wedge}\limits_{k=1}^{N} I_k\right),

  $$ 

где $N$~--- разрядность счетчика. На рис.~10 она описана рекурсивно и~при 

реализации допускает любую декомпозицию. 



\begin{figure} %fig10

\vspace*{1pt}

  \begin{center}

 \mbox{%

 \epsfxsize=125mm  %128mm

\epsfbox{ste-10.eps}

 }

\end{center}

\vspace*{-9pt}

\Caption{Параметризованное Verilog-описание модуля СС-счетчика со сбросом}

\vspace*{-6pt}

\end{figure}

  

  При автоматизированном синтезе СС-счет\-чи\-ка обязательно 

соблюдение по\-ляр\-ности сигналов управ\-ле\-ния, в~данном случае~--- активный 

передний фронт счетного входа~ТI и~активный низ\-кий уровень сигнала на 

входе начального сброса~Res. При не\-об\-хо\-ди\-мости по\-ляр\-ность входов 

иерархического модуля согласуется с~переменными в~исходном  

Verilog-опи\-са\-нии с~по\-мощью дополнительных инверторов. Кроме того, 

появляется инверсная компонента бифазного информационного выхода 

StateB. В~результате блок always с~асинхронным сбросом замещается 

фраг\-мен\-том Verilog-кода, показанным на рис.~11. 

  



\begin{figure} %fig11

\vspace*{1pt}

  \begin{center}

 \mbox{%

 \epsfxsize=125mm  %128mm

\epsfbox{ste-11.eps}

 }

\end{center}

\vspace*{-9pt}

\Caption{Подстановка модуля СС-счет\-чи\-ка со сбросом в~Verilog-опи\-сание }

\vspace*{-3pt}

\end{figure}

  

  Синтез СС-счетчика реализуется с~помощью следующей 

формализованной процедуры на примере исходного описания на рис.~1:

  \begin{enumerate}[(1)]

  \item \textbf{анализ заголовка блока always.} Цель~--- определить чис\-ло 

и~типы входов управ\-ле\-ния счетчиком, их ак\-тив\-ные уров\-ни и~фронты. 

В~данном случае заголовок содержит в~списке чув\-ст\-ви\-тель\-ности два 

сигнала: Т и~res. У~них обоих активным считается пе\-ред\-ний фронт~--- 

переключение сигнала из логического нуля в~логическую единицу. 

Следовательно, мож\-но предположить, что по крайней мере один из них 

служит сигналом синхронизации устройства, описываемого блоком always;

\item \textbf{специфицирование сигналов в~блоке always.} В~списке 

чув\-ст\-ви\-тель\-ности счетчика указываются счет\-ный сигнал (глобальный 

тактовый или производный от него), сигналы ини\-ци\-а\-ли\-за\-ции и~их 

активные уров\-ни или перепады. На рис.~1 таких сигналов два: Т и~res. 

Анализ тела блока always на рис.~1 показывает, что сигнал~Т в~теле 

блока не используется, сигнал же res высоким уровнем разрешает 

обнуление некоторой переменной State, а~низ\-ким уровнем разрешает 

инкремент этой же переменной. Значит, Т~--- синхросигнал с~активным 

передним фронтом, а~res~--- сигнал сброса. Просмотр объявлений 

переменных в~исходном Verilog-опи\-са\-нии всей схемы поз\-во\-ля\-ет 

определить раз\-ряд\-ность переменной State. Поскольку она 

инкрементируется в~блоке always, то это со\-сто\-яние счетчика, а~Т~--- его 

счетный вход с~активным передним фронтом; 

\item \textbf{подбор СС-аналога для синхронного устройства, 

описываемого данным блоком always.} Поскольку тело блока always 

описывает сум\-ми\-ру\-ющий счетчик, пе\-ре\-клю\-ча\-ющий\-ся по переднему 

фронту сигнала~Т и~об\-ну\-ля\-ющий\-ся по сигналу $\mathrm{res}=1$,  

в~СС-ва\-ри\-ан\-те он реализуется на счетном триггере~С1С из 

биб\-лио\-те\-ки~[21], к~входу сброса которого подключается инверсия 

сигнала res;

\item \textbf{cинтез индикаторной подсхемы.} Обязательный компонент 

индикаторной подсхемы СС-счет\-чи\-ка~--- индикатор счетного режима. 

Его целесообразно включать непосредственно в~модуль самого счетчика. 

Под\-тверж\-де\-ние успешной СС-предуста\-нов\-ки счетчика целесообразно 

формировать в~окру\-же\-нии счетчика;

\item \textbf{оформление счетчика в~виде отдельного модуля.} Блок 

always в~исходном Verilog-опи\-са\-нии син\-те\-зи\-ру\-емой схемы замещается 

иерархическим модулем СС-счет\-чи\-ка и~атрибутами, 

спе\-ци\-фи\-ци\-ру\-ющи\-ми его особенности функционирования. При 

по\-сле\-ду\-ющем синтезе атрибуты помогут корректно реализовать 

и~встроить СС-счет\-чик в~общую схему.

\end{enumerate}



\section{Заключение}

  

  Известные логические синтезаторы синхронных схем производят 

двоичные счетчики с~разными типами сброса и~установки единообразно~--- 

с~по\-мощью ре\-гист\-ра хранения с~параллельной записью следующего 

со\-сто\-яния счетчика, которое вы\-чис\-ля\-ет\-ся комбинационной схемой на основе 

текущего со\-сто\-яния и~сигнала, опре\-де\-ля\-юще\-го на\-прав\-ле\-ние счета. Все 

разряды такого счетчика об\-нов\-ля\-ют\-ся одновременно, что соответствует 

параллельной архитектуре счетчика. При наличии сброса или уста\-нов\-ки 

счетчика входы предуста\-нов\-ки формируются той же комбинационной 

схемой, что пред\-вы\-чис\-ля\-ет и~новое со\-сто\-яние счет\-чика. 

  

  Однако последовательный СС-счетчик обладает в~2--2,5~раза меньшей 

слож\-ностью и~в~1,7--2~раза луч\-шим быст\-ро\-дей\-ст\-ви\-ем в~сравнении 

с~параллельным СС-счет\-чи\-ком, поэтому для синтеза СС-счет\-чи\-ка 

предпочтительней использовать его по\-сле\-до\-ва\-тель\-ную реализацию.

  

  Самосинхронные счетчики целесообразно синтезировать с~использованием 

па\-ра\-мет\-ри\-зо\-ван\-ных шаб\-ло\-нов модуля счетчика и~его окружения, 

спе\-ци\-фи\-ци\-ру\-ющих особенности функционирования СС-счет\-чи\-ка 

с~различными вариантами предуста\-нов\-ки и~на\-прав\-ле\-ния счета. При этом 

модуль счетчика индицирует только свою работу в~счет\-ном режиме.  

Самосинхронная предуста\-нов\-ка и~на\-прав\-ле\-ние счета управ\-ля\-ют\-ся и~индицируются 

окружением СС-счет\-чи\-ка, которое организует необходимую процедуру  

СС-предуста\-нов\-ки или изменения на\-прав\-ле\-ния \mbox{счета.}



{\small\frenchspacing

{%\baselineskip=10.8pt

\addcontentsline{toc}{section}{Литература}

\begin{thebibliography}{99}

\bibitem{1-step}

\Au{Hennessy J.\,L., Patterson D.\,A.} Computer architecture: A~quantitative approach.~--- 6th 

ed.~--- San Mateo, CA, USA: Morgan Kaufmann, 2019. 936~p. 

\bibitem{2-step}

\Au{Spars{\o}~J.} Introduction to asynchronous circuit design.~--- Copenhagen, Denmark: 

DTU Compute, Technical University of Denmark, 2020. 255~p. {\sf 

https://backend.\linebreak orbit.dtu.dk/ws/files/215895041/JSPA\_async\_book\_2020\_PDF.pdf}.



\bibitem{4-step} %3

\Au{Jiang~W., Sha~E.\,H.-M., Zhuge~Q., Yang~L., Chen~X., Hu~J.} On the design of  

time-constrained and buffer-optimal self-timed pipelines~/\!\!/ IEEE T. Comput.

Aid.~D., 2019. Vol.~38. No.\,8. P.~1515--1528. doi: 

10.1109\!/\!TCAD.2018.2846642.

\bibitem{3-step} %4

\Au{Zakharov~V., Stepchenkov~Y., Diachenko~Y., Rogdestvenski~Y.} Self-timed circuitry 

retrospective~/\!\!/ Conference (International) on Engineering Technologies and Computer 

Science Proceedings.~--- Piscataway, NJ, USA: IEEE, 2020. P.~63--69. doi: 

10.1109\!/ EnT48576.2020.00018.

\bibitem{5-step}

\Au{Kushnerov A., Medina~M., Yakovlev~A.} Towards hazard-free multiplexer based 

implementation of self-timed circuits~/\!\!/ 27th Symposium (International) on 

Asynchronous Circuits and Systems Proceedings.~--- Piscataway, NJ, USA: IEEE, 2021. P.~17--24. doi: 

10.1109\!/\!ASYNC48570.2021.00011.





\bibitem{7-step} %6

\Au{Stepchenkov~Y.\,A., Kamenskih~A.\,N., Diachenko~Y.\,G., Rogdestvenski~Y.\,V., 

Diachenko~D.\,Y.} Improvement of the natural self-timed circuit tolerance to short-term soft 

errors~/\!\!/ Advances Science Technology Engineering Systems~J., 2020. Vol.~5. No.\,2.  

P.~44--56. doi: 10.25046\!/\!aj050206.



\bibitem{6-step} %7

   \Au{Соколов И.\,А., Степченков~Ю.\,А., Рождественский~Ю.\,В., Дьяченко~Ю.\,Г.} Приближенная 

оценка эффективности синхронной и~самосинхронной методологий в~задачах проектирования 

сбоеустойчивых вы\-чис\-ли\-тель\-но-управ\-ля\-ющих сис\-тем~/\!\!/ Автоматика и~телемеханика, 2022. №\,2.  

С.~122--132. doi: 10.31857\!/ S0005231022020088. EDN: PYRRXN. 



\bibitem{8-step}

\Au{Сурков А.\,В.} Использование Synopsys Design Compiler для синтеза 

самосинхронных схем~/\!\!/ Программные продукты и~системы, 2014. №\,4. С.~24--30. 

doi: 10.15827\!/ 0236-235X.108.024-030.

\bibitem{9-step}

\Au{Власов А.\,О., Сурков~А.\,В.} Маршрут проектирования самосинхронных 

конвейерных схем с~использованием возможностей САПР~/\!\!/ Программные 

продукты и~системы, 2015. №\,4. С.~110--115. doi: 10.15827\!/\!0236-235x.112.110-115.

\bibitem{10-step}

\Au{Edwards D., Bardsley~A., Jani~L., Plana~L., Toms~W.} Balsa: A~tutorial guide.~--- 

Manchester, U.K.: The University of Manchester, 2006. 165~p.  {\sf 

https://apt.cs.\linebreak manchester.ac.uk/ftp/pub/apt/balsa/3.5/BalsaManual3.5.pdf}.

\bibitem{11-step}

\Au{Spars{\o}~J., Furber~S.} Principles of asynchronous circuit design: A~systems 

perspective.~--- Dordrecht, Netherlands: Kluwer Academic Publs., 2001. 337~p. doi: 

10.1007\!/\!978-1-4757-3385-3.

\bibitem{12-step}

\Au{Andrikos~N., Lavagno~L., Pandini~D., Sotiriou~C.\,P.} A~fully automated 

desynchronization flow for synchronous circuits~/\!\!/ 44th ACM\!/\!IEEE Design Automation 

Conference Proceedings.~--- Piscataway, NJ, USA: IEEE, 2007. P.~982--985. doi: 

10.1145\!/\!1278480.1278722.

\bibitem{13-step}

\Au{Zhou R., Chong~K.-S., Gwee~B.-H., Chang~J.\,S., Ho~W.-G.} Synthesis of asynchronous 

QDI circuits using synchronous coding specifications~/\!\!/ Symposium (International) 

on Circuits and Systems Proceedings.~--- Piscataway, NJ, USA: IEEE, 2014. P.~153--156. doi: 

10.1109\!/\!ISCAS.2014.6865088.

\bibitem{14-step}

   \Au{Степченков Ю.\,А., Хилько~Д.\,В., Дьяченко~Ю.\,Г., Морозов~Н.\,В., 

Степченков~Д.\,Ю., Орлов~Г.\,А.} Методика десинхронизации при синтезе 

самосинхронных схем~/\!\!/ Системы и~средства информатики, 2024. Т.~34. №\,1.  

С.~33--43. doi: 10.14357\!/\!08696527240103. EDN: XGZCWU.

\bibitem{15-step}

\Au{Плеханов Л.\,П.} Основы самосинхронных электронных схем.~--- М.: Бином; 

Лаборатория знаний, 2013. 208~с.

\bibitem{16-step}

   \Au{Гаврилов С.\,В., Денисов~А.\,Н., Малашевич~Н.\,И., Росляков~А.\,С., 

Фёдоров~Р.\,А.} Диагностирование самосинхронных функциональных ячеек средствами 

САПР <<КОВЧЕГ>>~/\!\!/ Известия высших учебных заведений. Электроника, 2011. 

№\,1(87). С.~40--45. EDN: NDDTDX.

\bibitem{17-step}

\Au{Рождественский Ю.\,В., Морозов~Н.\,В., Рождественскене~А.\,В.} Под\-сис\-те\-ма 

событийного анализа самосинхронных схем АСПЕКТ~/\!\!/ Проб\-ле\-мы разработки 

перспективных микро- и~наноэлектронных сис\-тем.~--- М.: ИППМ РАН, 2010.  

С.~26--31.

\bibitem{18-step}

\Au{Зацаринный А.\,А., Степченков~Ю.\,А., Дьяченко~Ю.\,Г., Морозов~Н.\,В., 

Степченков~Д.\,Ю.} Автоматизация синтеза самосинхронных схем~/\!\!/ Сис\-те\-мы 

высокой доступности, 2023. Т.~19. №\,3. С.~48--56. doi:  

10.18127\!/\!j20729472-202303-04. EDN: IAPVXR.

\bibitem{19-step}

Yosys open synthesis suite. {\sf https://yosyshq.net/yosys}.

\bibitem{20-step}

   \Au{Степченков Ю.\,А., Степченков~Д.\,Ю., Дьяченко~Ю.\,Г., Морозов~Н.\,В., 

Плеханов~Л.\,П.}  Замена синхронных триггеров самосинхронными аналогами в~процессе 

десинхронизации схемы~/\!\!/ Сис\-те\-мы и~средства информатики, 2023. Т.~33. №\,4. 

С.~4--15. doi: 10.14357\!/\!08696527230401.  EDN: VPLSHI.



\bibitem{21-step}

\Au{Степченков Ю.\,А., Денисов~А.\,Н., Дьяченко~Ю.\,Г. и~др.} Библиотека 

функциональных ячеек для проектирования самосинхронных полузаказных БМК 

микросхем серий 5503\!/\!5507.~--- М.: Техносфера, 2017. 367~с. {\sf 

http://www.\linebreak technosphera.ru/lib/book/497}.

     \end{thebibliography}

} }







\vspace*{-6pt}



\hfill{\small\textit{Поступила в~редакцию 15.03.24}}



\vspace*{8pt}



\hrule



\vspace*{2pt}





\hrule



%\vspace*{8pt}



%\newpage



%\vspace*{-28pt}



\def\tit{SELF-TIMED COUNTER SYNTHESIS FORMALIZATION}







\def\titkol{Self-timed counter synthesis formalization}



\def\aut{Yu.\,A.~Stepchenkov, Yu.\,G.~Diachenko, N.\,V.~Morozov, D.\,Yu.~Stepchenkov, 

and~D.\,Yu.~Diachenko}



\def\autkol{Yu.\,A.~Stepchenkov, Yu.\,G.~Diachenko, N.\,V.~Morozov, et al.}

%D.\,Yu.~Stepchenkov,  and~D.\,Yu.~Diachenko}



\titel{\tit}{\aut}{\autkol}{\titkol}



\vspace*{-8pt}







\noindent

Federal Research Center ``Computer Science and Control'' of the Russian Academy of 

Sciences, 44-2~Vavilov Str., Moscow 119333, Russian Federation



\def\leftfootline{\small{\textbf{\thepage}

\hfill Sistemy i Sredstva Informatiki~---

Systems and Means of Informatics 2024 vol~34 no\,2}

}%

 \def\rightfootline{\small{Sistemy i Sredstva Informatiki~---

 Systems and Means of Informatics 2024 vol~34 no\,2

\hfill \textbf{\thepage}}}



\vspace*{3pt}



   

  

   

   \Abste{Self-timed (ST) circuits have high reliability. They guarantee detection and 

   localization of any persistent faults and demonstrate a~high level of 

   fault tolerance. However, designing ST circuits is more labor-intensive compared to synchronous circuits because 

   one should construct an additional indication subcircuit and adhere to the principles of truly ST circuit 

   implementation. Formalized desynchronization provides automatic conversion of the original 

   synchronous circuit description into the self-timed one but when synthesizing sequential ST units, 

   including ST counters, it leads to excessive hardware redundancy and, as a~consequence, to their low performance. 

   The article substantiates the approach to the ST counter synthesis based on the heuristic method formalization

    for their construction and ensuring the guaranteed resulted truly ST 

   implementation that functions in full accordance with the original description and has close to optimal consumer characteristics}

   

   \KWE{automated synthesis; self-timed circuit; counter; desynchronization; preset; indication}

   

\DOI{10.14357\!/\!08696527240205}{KDIEOJ}  



\vspace*{-10pt}





%\Ack

%\vspace*{-3pt}

%\noindent







%\vspace*{-3pt}



\renewcommand{\bibname}{\protect\rmfamily References}

%\renewcommand{\bibname}{\large\protect\rm References}



{\small\frenchspacing

{ %\baselineskip=12pt

\addcontentsline{toc}{section}{References}

\begin{thebibliography}{99}

   \bibitem{1-step-1}

   \Aue{Hennessy, J.\,L., and D.\,A.~Patterson.} 2019. \textit{Computer architecture: 

A~quantitative approach}. 6th ed. San Mateo, CA: Morgan Kaufmann. 936~p.

   \bibitem{2-step-1}

   \Aue{Spars{\o}, J.} 2020. \textit{Introduction to asynchronous circuit design}. Copenhagen, 

Denmark: DTU Compute, Technical University of Denmark. 255~p. 

Available at: {\sf 

https://backend.orbit.dtu.dk/ws/files/215895041/JSPA\_async\_book\_2020\_PDF.pdf} (accessed May~1, 2024).



   \bibitem{4-step-1} %3

   \Aue{Jiang, W., E.\,H.-M.~Sha, Q.~Zhuge, L.~Yang, X.~Chen, and J.~Hu.} 2019. On the design of  

time-constrained and buffer-optimal self-timed pipelines. \textit{IEEE T. Comput. Aid. D.} 

38(8):1515--1528. doi: 10.1109\!/\!TCAD.2018.2846642.



   \bibitem{3-step-1} %4

   \Aue{Zakharov, V., Y.~Stepchenkov, Y.~Diachenko, and Y.~Rogdestvenski.} 

 2020. Self-timed circuitry retrospective. \textit{Conference (International) on Engineering 

Technologies and Computer Science Proceedings}. Piscataway, NJ: IEEE. 63--69. doi: 10.1109\!/\!EnT48576.2020.00018.



   \bibitem{5-step-1}

   \Aue{Kushnerov, A., M.~Medina, and A.~Yakovlev.} 2021. Towards hazard-free 

multiplexer based implementation of self-timed circuits. \textit{27th Symposium 

(International) on Asynchronous Circuits and Systems Proceedings}. Piscataway, NJ: IEEE. 17--24. doi: 

10.1109\!/\!ASYNC48570.2021.00011.

   

   \bibitem{7-step-1} %6

   \Aue{Stepchenkov, Y.\,A., A.\,N.~Kamenskih, Y.\,G.~Diachenko, Y.\,V.~Rogdestvenski, 

and D.\,Y.~Diachenko.} 2020. Improvement of the natural self-timed circuit tolerance to  

short-term soft errors. \textit{Advances Science Technology Engineering Systems~J.} 

 5(2):44--56. doi: 10.25046\!/\!aj050206.

 

 \bibitem{6-step-1} %7

   \Aue{Sokolov, I.\,A., Yu.\,A.~Stepchenkov, Yu.\,V.~Rogdestvenski, and 

Yu.\,G.~Diachenko}. 2022. Approximate evaluation of the effectiveness of synchronous and 

self-timed methodologies in problems of designing failure-tolerant computing and control 

systems. \textit{Automat. Rem. Contr.} 83(2):264--272. doi: 10.1134\!/\!S0005117922020084.



   \bibitem{8-step-1}

   \Aue{Surkov, A.\,V.} 2014. Is\-pol'\-zo\-va\-nie Synopsys Design Compiler dlya sin\-te\-za 

sa\-mo\-sin\-khron\-nykh skhem [Synthesis of burst-mode asynchronous schemes using Synopsys 

Design Compiler]. \textit{Programmnye produkty i~sis\-te\-my} [Software \& Systems]  

4:24--30. doi: 10.15827\!/\!0236-235X.108.024-030.

   \bibitem{9-step-1}

   \Aue{Vlasov, A.\,O., and A.\,V.~Surkov.} 2015. Marsh\-rut pro\-ek\-ti\-ro\-va\-niya 

sa\-mo\-sin\-khron\-nykh kon\-veyer\-nykh skhem s~is\-pol'\-zo\-va\-ni\-em voz\-mozh\-no\-stey SAPR [Automated 

design flow of self-timed pipelines using EDA tools]. \textit{Programmnye produkty  

i~sis\-te\-my} [Software \&~Systems] 4:110--115. doi: 10.15827\!/\!0236-235x.112.110-115.

   \bibitem{10-step-1}

   \Aue{Edwards, D., A.~Bardsley, L.~Jani, L.~Plana, and W.~Toms.} 2006. Balsa: A~tutorial 

guide.  Manchester, U.K.: The University of Manchester. 165~p. Available at:  {\sf 

https://apt.cs.manchester.ac.uk/ftp/pub/apt/balsa/3.5/BalsaManual3.5.pdf} (accessed April~10, 2024).

   \bibitem{11-step-1}

   \Aue{Spars{\o},~J., and S.~Furber.} 2001. \textit{Principles of asynchronous circuit design: 

A~systems perspective}. Dordrecht, Netherlands: Kluwer Academic Publs. 337~p. doi: 

 10.1007\!/ 978-1-4757-3385-3.

   \bibitem{12-step-1}

   \Aue{Andrikos, N., L.~Lavagno, D.~Pandini, and C.\,P.~Sotiriou.} 2007. A~fully-automated 

desynchronization flow for synchronous circuits. \textit{44th ACM\!/\!IEEE Design Automation 

Conference Proceedings}. Piscataway, NJ: IEEE. 982--985. doi: 10.1145\!/ 1278480.1278722.

   \bibitem{13-step-1}

   \Aue{Zhou, R., K.-S.~Chong, B.-H.~Gwee, J.\,S.~Chang, and W.-G.~Ho.} 2014. Synthesis 

of asynchronous QDI circuits using synchronous coding specifications. \textit{Symposium 

(International) on Circuits and Systems Proceedings}. Piscataway, NJ: IEEE. 153--156. doi: 

10.1109\!/\!ISCAS.2014.6865088.

   \bibitem{14-step-1}

   \Aue{Stepchenkov, Yu.\,A., D.\,V.~Khilko, Yu.\,G.~Dia\-chen\-ko, N.\,V.~Mo\-ro\-zov, 

D.\,Yu.~Step\-chen\-kov, and G.\,A.~Or\-lov.} 2024. Me\-to\-di\-ka de\-sin\-khro\-ni\-za\-tsii pri sin\-te\-ze 

sa\-mo\-sin\-khron\-nykh skhem [Desynchronization methodology at self-timed circuit synthesis]. 

\textit{Sistemy i~Sredstva Informatiki~--- Systems and Means of Informatics} 34(1):33--43. doi: 

10.14357\!/\!08696527240103. EDN: XGZCWU.

   \bibitem{15-step-1}

   \Aue{Plekhanov, L.\,P.} 2013. \textit{Osno\-vy sa\-mo\-sin\-khron\-nykh elekt\-ron\-nykh skhem} 

[Basics of self-timed electronic circuits]. Moscow: Binom. Laboratoriya znaniy. 208~p.

   \bibitem{16-step-1}

   \Aue{Gavrilov, S.\,V., A.\,N.~Denisov, N.\,I.~Malashevich, A.\,S.~Ros\-lya\-kov, and 

R.\,A.~Fe\-do\-rov.} 2011. Diagno\-sti\-ro\-va\-nie sa\-mo\-sin\-khron\-nykh funk\-tsi\-o\-nal'\-nykh yache\-ek 

sredst\-va\-mi SAPR ``KOVCHEG'' [Diagnostics of self-synchronous functional cells by means of 

CAD ``Kovcheg'']. \textit{Izvestiya vysshikh uchebnykh zavedeniy. Elektronika} [Proceedings of 

Universities. Electronics] 1(87):40--45. EDN: NDDTDX.

   \bibitem{17-step-1}

   \Aue{Rogdestvenski, Y.\,V., N.\,V.~Morozov, and A.\,V.~Rog\-dest\-ven\-ske\-ne}. 2010. 

Pod\-sis\-te\-ma so\-by\-tiy\-no\-go ana\-li\-za sa\-mo\-sin\-khron\-nyh skhem ASPEKT [ASPECT: A~suite of  

self-timed event-driven analysis]. \textit{Problemy razrabotki perspektivnykh mikro- 

i~nanoelektronnykh sistem} [Problems of the perspective micro- and nanoelectronic systems 

development]. Moscow: IPPM RAN. 26--31.

   \bibitem{18-step-1}

   \Aue{Zatsarinnyy, A.\,A., Yu.\,A.~Stepchenkov, Yu.\,G.~Diachenko, N.\,V.~Morozov, and 

D.\,Yu.~Stepchenkov.} 2023. Av\-to\-ma\-ti\-za\-tsiya sin\-te\-za samosinkhronnykh skhem [Self-timed 

circuits automated design]. \textit{Sistemy vysokoy dostupnosti} [Highly Available Systems] 

19(3):48--56. doi: 10.18127\!/\!j20729472-202303-04. EDN: IAPVXR.

   \bibitem{19-step-1}

Yosys open synthesis suite. Available at: {\sf https://yosyshq.net/yosys} (accessed April~10, 2024).

   \bibitem{20-step-1}

   \Aue{Stepchenkov, Yu.\,A., D.\,Yu.~Stepchenkov, Yu.\,G.~Diachenko, N.\,V.~Morozov, 

and L.\,P.~Plekhanov.} 2023. Za\-me\-na sinkhronnykh trig\-ge\-rov samosinkhronnymi ana\-lo\-ga\-mi 

v~pro\-tses\-se desinkhronizatsii skhe\-my [Replacing synchronous triggers with self-timed 

counterparts during circuit desynchronization]. \textit{Sistemy i~Sredstva Informatiki~--- 

Systems and Means of Informatics} 33(4):4--15. doi: 10.14357\!/\!08696527230401. EDN: 

VPLSHI.

   \bibitem{21-step-1}

   \Aue{Stepchenkov, Yu.\,A., A.\,N.~Denisov, Yu.\,G.~Diachenko, \textit{et al.}} 2017. 

\textit{Bib\-lio\-te\-ka funk\-tsi\-o\-nal'\-nykh yache\-ek dlya pro\-ek\-ti\-ro\-va\-niya sa\-mo\-sin\-khron\-nykh po\-lu\-za\-kaz\-nykh 

BMK mik\-ro\-skhem se\-riy 5503\!/\!5507} [Functional cell library for designing self-timed  

semi-custom integrated circuits on 5503\!/\!5507 gate arrays]. Moscow: Tekhnosfera. 367~p. 

Available at: {\sf http://www.technosphera.ru/lib/book/497} (accessed April~10, 2024). 

   

   

 \end{thebibliography}

} }



\vspace*{-4pt}



\hfill{\small\textit{Received March 15, 2024}} 



\vspace*{-7pt} 



\Contr



\vspace*{-3pt}

   

   \noindent

   \textbf{Stepchenkov Yuri A.} (b.\ 1951)~--- Candidate of Science (PhD) in technology, head 

of department, leading scientist, Federal Research Center ``Computer Science and Control'' of 

the Russian Academy of Sciences, 44-2~Vavilov Str., Moscow 119133, Russian Federation; 

\mbox{YStepchenkov@ipiran.ru}



\pagebreak

   

%   \vspace*{3pt}

   

   \noindent

   \textbf{Diachenko Yuri G.} (b.\ 1958)~--- Candidate of Science (PhD) in technology, senior 

scientist, Federal Research Center ``Computer Science and Control'' of the Russian Academy of 

Sciences, 44-2~Vavilov Str., Moscow 119133, Russian Federation; \mbox{diaura@mail.ru}

   

   \vspace*{3pt}

   

   \noindent

   \textbf{Morozov Nikolai V.} (b.\ 1956)~--- senior scientist, Federal Research Center 

``Computer Science and Control'' of the Russian Academy of Sciences, 44-2~Vavilov Str., 

Moscow 119133, Russian Federation; \mbox{NMorozov@ipiran.ru}

   

   \vspace*{3pt}

   

   \noindent

   \textbf{Stepchenkov Dmitri Yu.} (b.\ 1973)~--- senior scientist, Federal Research Center 

``Computer Science and Control'' of the Russian Academy of Sciences, 44-2~Vavilov Str., 

Moscow 119133, Russian Federation; \mbox{stepchenkov@mail.ru}

   

   \vspace*{3pt}

   

   \noindent

   \textbf{Diachenko Denis Yu.} (b.\ 1987)~--- engineer-researcher, Federal Research Center 

``Computer Science and Control'' of the Russian Academy of Sciences, 44-2~Vavilov Str., 

Moscow 119133, Russian Federation; \mbox{diaden87@gmail.com}



\label{end\stat}



\renewcommand{\bibname}{\protect\rm Литература}

